El \textit{pipeline} de ingesta y reconocimiento convierte cada fichero
escaneado que llega a la carpeta \acs{sftp} en un \texttt{ExamPage}
enriquecido con datos estructurados, y asigna el conjunto resultante a
las \dfnpl{instancia} de los alumnos convocados. Se compone de tres
subtareas Celery diferenciadas que se ejecutan en cascada:
\texttt{ingest\_page\_task} (normalización),
\texttt{recognize\_page\_task} (reconocimiento) y
\texttt{assemble\_exam\_pages} (ensamblado y \textit{matching},
diferida). Las tres se diagraman por separado en las
\cref{fig:ingest-page-task,fig:recognize-page-task,fig:assemble-exam-pages};
la composición completa de las tres en su orden de ejecución natural se
incluye en la \cref{fig:ocr-flow-full} (\cref{AP:SEQUENCE}).


\subsection{Normalización}
\label{SS:DES-INGEST-PROCESSING}

La tarea \texttt{ingest\_page\_task} es el punto de entrada del
\textit{pipeline} (\cref{fig:ingest-page-task}). Recibe un fichero
—imagen o \acs{pdf} multipágina— depositado por el \dfn{watcher} de
\acs{sftp}, valida su tamaño contra el límite configurado por la
organización y lo normaliza a una o varias imágenes png a
\(200\,\text{dpi}\). Cada imagen se sube a MinIO bajo una clave que
codifica organización, examen y página, y se crea un registro
\texttt{ExamPage} en estado \texttt{PENDING\_RECOGNITION}. Finalmente, se
encola una tarea \texttt{recognize\_page\_task} por página en la cola
\texttt{recognition} de Celery.

\begin{figure}[Subtarea de normalización]%
              {fig:ingest-page-task}%
              {Subtarea \texttt{ingest\_page\_task}: normalización del
               fichero entrante y creación de las \texttt{ExamPage}.}
  \image{}{0.32\textheight}{plantuml/ingest_page_task}
\end{figure}

La conservación del lote físico de origen es relevante para la fase
posterior de \textit{matching}: las páginas extraídas de un mismo
fichero comparten un \texttt{batch\_id} que las identifica como
pertenecientes al mismo examen, lo que aporta evidencia cruzada cuando
alguna de ellas tiene un reconocimiento débil.


\subsection{Reconocimiento}
\label{SS:DES-INGEST-RECOGNITION}

La tarea \texttt{recognize\_page\_task} (\cref{fig:recognize-page-task})
descarga la imagen de MinIO e invoca al \textit{dispatcher} de
reconocedores. El reconocimiento procede en dos pasos:

\begin{enumerate}
  \item \textbf{Decodificación \acs{qr}.} Un único \texttt{QRRecognizer}
  se ejecuta sobre la imagen completa para identificar la página. Si el
  \acs{qr} decodifica correctamente y supera la verificación de
  \textit{checksum}, se extraen \texttt{org\_id}, \texttt{exam\_id},
  \texttt{model\_id} y \texttt{page\_number}, y se resuelve el
  \texttt{PageProfile} correspondiente, que define las
  \dfnpl{zonaocr} esperadas en esa página. Si el \acs{qr} es ilegible, la
  página se marca con incidencia \texttt{QR\_READ\_ERROR} y se omite el
  resto del reconocimiento: la página permanece en
  \texttt{PENDING\_RECOGNITION} a la espera de intervención manual del
  coordinador.

  \item \textbf{Reconocimiento de zonas.} Para cada zona no-\acs{qr} del
  perfil, el \textit{dispatcher} recorta la región correspondiente y
  delega en el reconocedor adecuado según el tipo de zona: el
  \texttt{OCRRecognizer} para texto y números —apoyado en el motor
  \acs{ocr} \dfn{plugin} configurado por la organización—, o el
  \texttt{CheckboxRecognizer} para casillas de verificación, basado en
  densidad de píxeles oscuros.
\end{enumerate}

\begin{figure}[Subtarea de reconocimiento]%
              {fig:recognize-page-task}%
              {Subtarea \texttt{recognize\_page\_task}: decodificación
               \acs{qr} y reconocimiento de zonas del perfil mediante el
               motor activo y los reconocedores especializados.}
  \image{}{0.40\textheight}{plantuml/recognize_page_task}
\end{figure}

Los resultados se almacenan en el campo \texttt{recognized\_data} del
\texttt{ExamPage}, junto con la confianza reportada por el motor para
cada zona, y la página transita a \texttt{RECOGNIZED}. A continuación se
dispara \texttt{schedule\_assembly}, que encola la subtarea de
ensamblado con una ventana de espera configurable para permitir que se
acumulen el resto de páginas del mismo examen antes de resolver el
\textit{matching} sobre el conjunto completo.


\subsection{Ensamblado y asignación óptima}
\label{SS:DES-INGEST-ASSEMBLY}

El problema del ensamblado consiste en decidir a qué alumno pertenece
cada página reconocida. Por la mecánica de la entrega física —el alumno
escribe nombre, \acs{nia} y \acs{dni} en zonas \acs{ocr} de la portada
y, opcionalmente, en páginas interiores— el sistema dispone de hasta
cuatro campos identificativos por examen, todos ellos contaminados por
errores de escritura y por errores de reconocimiento. El diseño actual
(V2) llega tras dos formulaciones intermedias que conviene enunciar
para entender qué problema resuelve cada salto.

\paragraph{V1 — Decisión local por página.}
La formulación inicial trataba el problema como una clasificación
independiente por página: para cada \texttt{ExamPage}, se calculaba la
similaridad de sus campos \acs{ocr} contra cada alumno convocado y se
le asignaba al alumno con mayor puntuación si esta superaba un umbral
mínimo. Esta formulación falla por construcción. En primer lugar, las
decisiones locales no garantizan biyección: dos páginas distintas
pueden caer en el mismo alumno y otro queda sin examen, sin que el
algoritmo tenga manera de resolver el conflicto. En segundo lugar, son
incoherentes entre páginas del mismo lote físico, que por construcción
pertenecen al mismo alumno: una página puede asignarse a un alumno y
otra del mismo \acs{pdf} a otro distinto, en función del ruido \acs{ocr}
que afecte a cada portada.

\paragraph{V1.5 — \textit{Greedy} con descartes y refuerzo por lote.}
La segunda formulación corrige los dos defectos anteriores agregando
las páginas en lotes (\texttt{batch\_id}) y resolviendo los lotes en
orden de confianza descendente, retirando del pool de candidatos a cada
alumno cuando se le asigna un lote. El refuerzo por lote estabiliza las
decisiones entre páginas hermanas —la evidencia \acs{ocr} de varias
páginas vota por el mismo candidato—, y el procesamiento ordenado
garantiza biyección. Sin embargo, sigue siendo una heurística
\textit{greedy} con el orden de procesamiento como única salvaguarda: un
lote con \textit{matching} débil que se procese antes que otro con
\textit{matching} fuerte puede consumir al mejor candidato del segundo,
sin que la formulación tenga manera de revisar la decisión a la luz de
la nueva evidencia. El resultado no es óptimo globalmente.

\paragraph{V2 — Optimización global mediante el algoritmo húngaro.}
La formulación adoptada redefine el problema como una asignación global
\(1\)-a-\(1\) entre lotes y alumnos convocados, equivalente al problema
clásico de asignación lineal: dada una matriz de puntuaciones \(P_{ij}\)
entre el lote~\(i\) y el alumno~\(j\), encontrar la asignación \(M\)
que maximice
\(\sum_{(i,j)\in M} P_{ij}\) sujeta a que cada lote y cada alumno
aparezcan en como mucho una pareja. Este problema se resuelve en tiempo
polinómico \(\Theta(n^{3})\) mediante el algoritmo húngaro de
Kuhn--Munkres~\cite{kuhn1955hungarian}, con la implementación de
\texttt{scipy.optimize.linear\_sum\_assignment}. La asignación así
obtenida es globalmente óptima, garantiza biyección por construcción y
procesa la información del examen completo a la vez en lugar de en
orden de llegada.

\begin{figure}[Subtarea de ensamblado y matching]%
              {fig:assemble-exam-pages}%
              {Subtarea \texttt{assemble\_exam\_pages}: ventana de
               espera, construcción de la matriz de puntuaciones,
               algoritmo húngaro y persistencia de las
               \texttt{ExamInstance}.}
  \image{}{0.36\textheight}{plantuml/assemble_exam_pages}
\end{figure}

La \cref{fig:assemble-exam-pages} muestra el flujo concreto de la
subtarea. Para evitar resolver el \textit{matching} antes de tiempo y
desperdiciar trabajo cuando aún faltan páginas por llegar, la subtarea
se dispara mediante una ventana de espera con doble criterio: por
\textbf{saturación}, cuando han llegado todas las páginas esperadas
\((\text{n\_alumnos}\times\text{páginas\_por\_modelo})\); o por
\textbf{timeout}, si vence el plazo de espera. La política optimista
sobre el umbral es deliberada: los lotes cuyo \textit{matching} no
supere el umbral mínimo se persisten igualmente como
\dfnpl{instancia} sin identificar marcadas con la incidencia
\texttt{LOW\_MATCH\_CONFIDENCE}, dejando la decisión final al
\dfn{corrector} en el flujo de revisión. Cuesta menos a un corrector
revisar un examen mal asignado que rastrear un examen ausente, y el
ciclo de corrección ya incluye una etapa de revisión humana en la que
cualquier asignación errónea aflora.

La V3 contemplada en la \cref{SEC:CONCL-FUTURE-WORK} extiende esta
formulación con refuerzo cruzado entre páginas del mismo lote
(\textit{cross-page reinforcement}), aprovechando que la caligrafía y
los errores \acs{ocr} suelen propagarse coherentemente entre páginas
escritas por la misma persona. El detalle de implementación de V2 —la
función de \textit{scoring}, las clases de equivalencia \acs{ocr} y la
construcción de la matriz de costes— se desarrolla en la
\cref{SS:IMPL-MATCHING}.
